% Данный файл распространяетсяя под лицензией CC BY 4.0. Текст лицензии размещён на https://creativecommons.org/licenses/by/4.0/
% (c) Кафедра системного программирования, 2025

\documentclass[a4paper]{article}

\usepackage[a4paper, top=8mm, bottom=8mm, left=8mm, right=8mm]{geometry}

\usepackage{polyglossia}
\setdefaultlanguage[babelshorthands=true]{russian}
\setotherlanguage{english}

\usepackage{fontspec}
\setmainfont{FreeSerif}
\newfontfamily{\russianfonttt}[Scale=0.7]{DejaVuSansMono}

\usepackage[tiny, compact]{titlesec}

\usepackage{titling}
\setlength{\droptitle}{-1cm}
\pretitle{\begin{center}\begin{bfseries}\Large}
\posttitle{\par\end{bfseries}\end{center}}
\preauthor{\begin{center}\normalsize}
\postauthor{\par\end{center}\vspace{-1.8cm}}

\usepackage{hyperref}
\usepackage{bookmark}
\usepackage{csquotes}
\usepackage{booktabs}
\usepackage{float}
\usepackage{enumitem}
\usepackage{amsmath}
% Сюда можно дописывать нужные вам пакеты.

\title{Реализация SAT-решателя с эвристикой Jeroslaw-Wang}

\author{Тупикова Екатерина Викторовна}

% Дата много места занимает, её лучше не указывать.
\date{}

\begin{document}

\maketitle

\begin{flushright}
    Группа: \emph{25.Б42-мм}

    Кафедра: \emph{системного программирования}

    % Степень/звание и должность научника мы лучше вас знаем, их можно не указывать.
    Научный руководитель: \emph{Смирнов Кирилл Константинович}

    Номер семестра практики: \emph{2}
\end{flushright} 


\section{Введение}
Задача выполнимости булевых формул (Propositional Satisfiability, SAT) заключается в определении, существует ли такое назначение истинностных значений переменным, при котором заданная булева формула принимает значение «истина».

Литерал — это переменная или её логическое отрицание.

Дизъюнкт представляется как дизъюнкция литералов, а конъюнктивно нормальной формой (КНФ) является конъюнкция конечного количества дизъюнктов. 

DIMACS — стандартный формат файлов для представления задач выполнимости. 

К задаче SAT сводятся задачи формальной верификации, автоматического
планирования, составления расписаний и решения комбинаторных задач. Поскольку
SAT является NP-полной задачей, практическая производительность полного решателя
во многом определяется порядком исследования возможных присваиваний.

\section{Постановка задачи}

Целью данной работы является реализация и экспериментальное исследование эвристики ветвления Jeroslow-Wang, описанной в статье \cite{marques1999}, а также оценка её влияния на производительность SAT-решателя на различных классах задач.

\section{Реализация}

Для реализации был выбран язык C++, так как он знаком учащейся.

Ядром SAT-солвера является DPLL (алгоритм Дэвиса — Патнема — Логемана — Лавленда) — алгоритм полного перебора, включающий себя принципы метода ветвей и границ, а так же использующий распространение единичных дизъюнктов.

На каждом шаге решатель выбирает свободную переменную, распространяет
вынужденные присваивания из единичных дизъюнктов и выполняет откат при
обнаружении конфликта. Если все дизъюнкты удовлетворены, найденное присваивание
является решением.

Без эвристики переменные рассматриваются в порядке их номеров:
\begin{equation}
    x = \min \{x_i \mid x_i \text{ ещё не назначена}\}.
\end{equation}
Такой способ прост и почти не требует дополнительных вычислений, однако его
эффективность зависит от порядка переменных во входном файле.

В качестве оптимизации процесса выбора переменных была выбрана эвристика Jeroslow-Wang. Её работа базируется на оценке встречаемости литералов в невыполненных дизъюнктах. Для заданного литерала $l$ вычисляется весовая функция $J(l)$ по следующей формуле:

\begin{equation}
    J(l) = \sum_{l \in \omega \wedge \omega \in \varphi} 2^{-|\omega|}
\end{equation}

где $\omega$~--- дизъюнкт, содержащий литерал $l$, а $|\omega|$~--- текущее количество свободных литералов в данном дизъюнкте.

Существует две вариации данной эвристики:
\begin{itemize}
    \item \textbf{Односторонняя эвристика Jeroslow-Wang (JW-OS)}: на каждом шаге выбирается присваивание, которое удовлетворяет литерал с наибольшим значением весовой функции $J(l)$.
    \item \textbf{Двусторонняя эвристика Jeroslow-Wang (JW-TS)}: алгоритм определяет переменную $x$ с максимальной суммарной оценкой $J(x) + J(\neg x)$.
\end{itemize}
Была выбрана двусторонняя вариация, поскольку она является более лёгкой в реализации.

В реализованной JW-TS сначала выбирается переменная
\begin{equation}
    x^* = \underset{x}{\operatorname{arg\,max}}\,
          \bigl(J(x) + J(\neg x)\bigr),
\end{equation}
после чего выбирается знак с большим индивидуальным весом:
\begin{equation}
    \nu(x^*) =
    \begin{cases}
        1, & J(x^*) \geq J(\neg x^*),\\
        0, & J(x^*) < J(\neg x^*).
    \end{cases}
\end{equation}

Множитель $2^{-|\omega|}$ задаёт больший приоритет коротким дизъюнктам, которые
находятся ближе к превращению в единичные или конфликтные. Веса динамически
обновляются при присваиваниях и откатах, что позволяет учитывать текущее
состояние формулы, но создаёт дополнительные вычислительные расходы.

\section{Эксперименты}

Для оценки эффективности эвристики проводились замеры времени работы на нескольких классах данных.

Эксперименты проводились на системе со следующими характеристиками:

\begin{itemize}
  \item Процессор: Intel Core i5-1235U
  \item ОЗУ: 16 ГБ
  \item ОС: Fedora Linux 43
\end{itemize}

Для каждого класса измерялось суммарное время последовательного решения всех
входящих в него формул. Сравнивались две конфигурации одного и того же
решателя: наивный выбор первой свободной переменной и выбор переменной с
помощью JW-TS.

Замеры выполнялись утилитой \texttt{hyperfine}: каждый тест повторялся 30 раз
после предварительного прогрева, после чего фиксировались среднее время,
стандартное отклонение и относительное ускорение.

Относительное ускорение JW-TS вычислялось как
\begin{equation}
    S = \frac{T_{\text{naive}}}{T_{\text{JW-TS}}},
\end{equation}
где $T_{\text{naive}}$ и $T_{\text{JW-TS}}$ --- средние времена решения всей
группы файлов. При $S>1$ применение эвристики ускоряет поиск, а при $S<1$
приводит к замедлению.

В качестве тестовых наборов были выбраны случайные выполнимые и невыполнимые
3-КНФ семейства UF и UUF, формулы кодирования задачи раскраски графа
\texttt{flat50-115}, а также невыполнимые формулы, кодирующие принцип Дирихле
(\texttt{pigeon-hole}). Наборы UF и UUF позволяют сравнить работу эвристик на
случайных формулах вблизи фазового перехода, тогда как два последних набора
имеют выраженную структуру и большое количество бинарных дизъюнктов.

UF и UUF представляют выполнимые и невыполнимые случайные 3-КНФ. Набор
\texttt{flat50-115} получен кодированием раскраски графа, а
\texttt{pigeon-hole} --- кодированием невыполнимого размещения большего числа
объектов в меньшем числе ячеек.

\begin{table}[H]
    \centering
    \caption{Сравнение производительности (время выполнения всей группы файлов, с)}
    \label{tab:results}
    \vspace{0.5em}
    \begin{tabular}{lccc}
        \toprule
        \textbf{Класс данных} & \textbf{Наивное ветвление} & \textbf{Эвристика JW-TS} & \textbf{Ускорение}\\
        \midrule
        uf50-218       & $1.274 \pm 0.008$ & $0.739 \pm 0.006$ & $1.72$\\
        uuf50-218      & $2.750 \pm 0.019$ & $2.266 \pm 0.018$ & $1.21$\\
        uf75-325       & $1.416 \pm 0.006$ & $0.359 \pm 0.006$ & $3.95$\\
        uuf75-325      & $2.945 \pm 0.007$ & $1.366 \pm 0.004$ & $2.16$\\
        flat50-115     & $1.389 \pm 0.019$ & $1.781 \pm 0.018$ & $0.78$\\
        pigeon-hole    & $0.989 \pm 0.024$ & $1.660 \pm 0.014$ & $0.60$\\
        \bottomrule
    \end{tabular}
\end{table}

На случайных 3-КНФ эвристика Jeroslow-Wang во всех случаях ускорила решение.
Наибольший выигрыш получен на выполнимом наборе \texttt{uf75-325}: время
сократилось почти в четыре раза. При увеличении числа переменных цена неудачного
раннего ветвления возрастает из-за быстрого роста дерева поиска, поэтому выбор
переменных по их весу начинает сильнее влиять на итоговое время. На
невыполнимых формулах ускорение оказалось ниже, чем на выполнимых формулах того
же размера. Для доказательства невыполнимости требуется исключить все
существенные ветви дерева, тогда как для выполнимой формулы достаточно найти
одно удовлетворяющее присваивание.

На структурированных наборах \texttt{flat50-115} и \texttt{pigeon-hole}
эвристика, напротив, увеличила время выполнения. В формулах раскраски графа
порядок переменных уже соответствует структуре исходной задачи, поэтому
последовательное наивное ветвление оказывается достаточно удачным. Формулы
\texttt{pigeon-hole} обладают высокой симметрией, из-за которой многие
переменные получают близкие веса. Кроме того, оба набора содержат преимущественно
бинарные дизъюнкты, обеспечивающие активное распространение единичных
ограничений даже без Jeroslow-Wang. В результате эвристика недостаточно
сокращает дерево поиска, а расходы на динамический пересчёт весов приводят к
замедлению. Таким образом, эффективность Jeroslow-Wang определяется не только
размером формулы, но и её структурой.

\subsection{Сопоставление с результатами Marques-Silva}

Полученные результаты согласуются с выводами исследования
Marques-Silva~\cite{marques1999}. В нём показано, что выбор эвристики ветвления
существенно влияет на отдельные классы и экземпляры SAT, однако ни одна из
рассмотренных эвристик не является лучшей для всех задач. В частности,
результаты автора для GRASP различаются между классами формул: JW-TS показывает
хорошее время на \texttt{aim-200}, но заметно уступает другим стратегиям на
\texttt{dubois}, \texttt{ssa} и ряде структурированных наборов. Аналогичная
зависимость наблюдается в проведённом эксперименте: JW-TS ускоряет решение
случайных 3-КНФ, но проигрывает на формулах раскраски графа и принципа Дирихле.

\subsection{Ограничения эксперимента}

Измерялось суммарное время групп формул, поэтому времена разных наборов нельзя
сравнивать без учёта числа файлов. Кроме того, решатель не использует обучение
на конфликтах и нехронологический откат, а число ветвлений и откатов отдельно
не измерялось. Для более полного исследования следует собирать эти показатели
и повторять тесты после случайной перенумерации переменных.

\section{Заключение}

В ходе работы был реализован SAT-решатель на основе алгоритма DPLL с
распространением единичных дизъюнктов. В решатель была добавлена двусторонняя
эвристика Jeroslow-Wang, динамически учитывающая текущие длины невыполненных
дизъюнктов. Для оценки её эффективности проведено сравнение с наивным
последовательным выбором переменных на случайных и структурированных наборах
формул.

Эксперименты показали, что Jeroslow-Wang особенно эффективна на случайных
3-КНФ: на исследованных наборах она обеспечила ускорение от 1.21 до 3.95 раза.
При увеличении размера формул преимущество эвристики возрастало, поскольку
удачный выбор переменных позволял исключать более крупные части дерева поиска.
В то же время на структурированных формулах раскраски графа и принципа Дирихле
эвристика привела к замедлению. В этих задачах последовательный порядок
переменных уже согласован со структурой ограничений, а дополнительные расходы
на пересчёт весов не компенсируются достаточным сокращением дерева поиска.

Таким образом, эвристика Jeroslow-Wang способна существенно повысить
производительность DPLL-решателя, однако её эффективность зависит от структуры
входной формулы. Полученные результаты подтверждают необходимость учитывать не
только качество выбора переменных, но и стоимость самой эвристики. Дальнейшее
развитие решателя может включать сбор статистики о количестве ветвлений и
откатов, а также реализацию более современных методов, таких как обучение на
конфликтах и адаптивный выбор эвристики.

\renewcommand{\refname}{Список литературы}
\begin{thebibliography}{9}
    \bibitem{marques1999}
    Marques-Silva J. The Impact of Branching Heuristics in Propositional Satisfiability Algorithms // EPIA'99, LNAI 1695. --- Springer-Verlag Berlin Heidelberg, 1999. --- P. 62-74.

    \bibitem{SATLIB}
    SATLIB: наборы задач для тестирования SAT-решателей [Электронный ресурс].
    URL: \url{https://www.cs.ubc.ca/~hoos/SATLIB/benchm.html}
    (дата обращения: 13.06.2026).

    \bibitem{DIMACS cnf}
    Описание формата DIMACS CNF [Электронный ресурс].
    URL: \url{https://people.sc.fsu.edu/~jburkardt/data/cnf/cnf.html}
    (дата обращения: 13.06.2026).

    \bibitem{repository}
    Репозиторий проекта «SAT-solver» [Электронный ресурс].
    URL: \url{https://github.com/sugahu565/SAT-solver}
    (дата обращения: 13.06.2026).
\end{thebibliography}


\end{document}
